\section{Introduction}

We consider the semilinear damped wave equation
\begin{equation}\label{eq}
\begin{cases}
\partial_{tt}^2 u + \gamma(x)\partial_t u = \Delta u - \alpha u - f(x,u)\:&(x,t)\in\Omega\times (0,+\infty)\\
u_{|\partial\Omega}(x,t)=0&(x,t)\in\partial\Omega\times (0,+\infty)\\
\left(u(\cdot,t=0),\partial_t u(\cdot,t=0)\right)=U_0=(u_0,u_1)& \mkern43mu\in H^1_0(\Omega)\times L^2(\Omega)
\end{cases}
\end{equation}
in the following general framework:
\begin{enumerate}\romanenumi
\item \label{listintro1.1}the domain $\Omega$ is a two-dimensional smooth compact and connected manifold with or without smooth boundary. If $\Omega$ is not flat, $\Delta$ has to be taken as Beltrami Laplacian operator.
\item \label{listintro1.2}the constant $\alpha\geq 0$ is a non-negative constant. We require that $\alpha>0$ in the case without boundary to ensure that $\Delta-\alpha Id$ is a negative definite self-adjoint operator.
\item \label{listintro1.3}the damping $\gamma\in L^\infty(\Omega,\RR_+)$ is a bounded function with non-negative values. Since we want to consider a damped equation, we will assume that $\gamma$ does not vanish everywhere.
\item \label{listintro1.4}the non-linearity $f\in\Cc^1(\overline\Omega\times\RR,\RR)$ is of polynomial type in the sense that there exists a constant $C$ and a power $p\geq 1$ such that for all $(x,u)\in \overline\Omega\times\RR$,
\begin{equation}\label{hyp-f-estim}
|f(x,u)|+|\grad_xf(x,u)|\leq C(1+|u|)^p\:\text{ and }\:|f'_u(x,u)|\leq C(1+|u|)^{p-1}\,.
\end{equation}
Moreover, in most of this paper, we will be interested in the stabilization problem and we will also assume that
\begin{equation}\label{hyp-f-sign}
\forall (x,u)\in \overline\Omega\times\RR\,,\:f(x,u)u\geq 0\,.
\end{equation}
\end{enumerate}


We introduce the space $X=H^1_0(\Omega)\times L^2(\Omega)$ and the operator $A$ defined by
\[
D(A)=\left(H^2(\Omega)\cap H^1_0(\Omega)\right)\times H^1_0(\Omega)\quad
A =
\begin{pmatrix}
0 & Id \\
\Delta-\alpha Id & -\gamma(x)
\end{pmatrix}\,.
\]
In this paper, we are interested in the cases where the linear semigroup $e^{At}$ has no uniform decay, that is that $\|e^{At}\|_{\Lc(X)}$ does no converge to zero. We only assume a semi-uniform decay, but sufficiently fast in the following sense.

\begingroup
\begin{enumerate}\romanenumi
\setcounter{enumi}{4}
\item \label{listintro1.5}There exists a function $h(t)$ such that
\begin{equation}\label{hyp-decay-1}
\forall U_0\in D(A)\,,\,\left\|e^{At}U_0\right\|_X \, \leq \, h(t)\,\|U_0\|_{D(A)}
\end{equation}
and there is $\sigma_h\in(0,1]$ such that
\begin{equation}\label{hyp-decay-2}
\lim_{t\rightarrow \infty}h(t)=0\;\text{ and }\,
\forall \sigma\in[0,\sigma_h]\,,\:\int_0^\infty h(t)^{1-\sigma}\,\drm t\,<\,\infty\,.
\end{equation}
\end{enumerate}
\endgroup
Condition~\eqref{hyp-decay-2} requires a decay rate fast enough to be integrable. Roughly speaking, this article shows that this condition, together with a suitable unique continuation property, are sufficient to obtain a stabilization of the semilinear equation. The relevant unique continuation property is explained in Proposition~\ref{prop-UCP-implique-th} below. We present two general results where it can be obtained.

Our first result concerns analytic nonlinearities and smooth dampings.

\begin{theo}\label{th1}
Consider the damped wave equation~\eqref{eq} in the framework of Assumptions~\eqref{listintro1.1}-\eqref{listintro1.5}. Assume in addition that:
\begin{enumerate}\alphenumi
\item \label{listtheo1.1.1}the function $(x,u)\mapsto f(x,u)$ is smooth and analytic with respect to $u$.
\item \label{listtheo1.1.2}the damping $\gamma$ is of class $\Cc^1$ or at least that there exists $\tilde\gamma\in\Cc^1(\overline\Omega,\RR_+)$ such that~\eqref{listintro1.5} holds with $\gamma$ replaced by $\tilde\gamma$ and such that the support of $\tilde\gamma$ is contained in the support of $\gamma$.
\item \label{listtheo1.1.3}the power $p$ of $f$ in~\eqref{hyp-f-estim} and the decay rate $h(t)$ of the semigroup in~\eqref{hyp-decay-1} satisfy $h(t)=\Oc(t^{-\beta})$ with $\beta>2p$.
\end{enumerate}
Then, any solution $u$ of~\eqref{eq} satisfies
\[
\|(u,\partial_t u)(t)\|_{H^1_0\times L^2}\,\xrightarrow[\:t\longrightarrow
+\infty\:]{}\,0\,.
\]
Moreover, for any $R$ and $\sigma>0$, there exists $h_{R,\sigma}(t)$ which goes to zero when $t$ goes to $+\infty$ such that the following stabilization hold. For any $U_0\in H^{1+\sigma}_0(\Omega)\times H^\sigma(\Omega)$, if $u$ is the solution of~\eqref{eq}, then
\[
\|(u_0,u_1)\|_{H^{1+\sigma}\times H^\sigma}\leq
R\:\Longrightarrow\:\|(u,\partial_t u)(t)\|_{H^1_0\times L^2} \leq
h_{R,\sigma}(t)\,\xrightarrow[\:t\longrightarrow +\infty\:]{}\,0\,.
\]
\end{theo}

Our assumptions~\eqref{listintro1.5} and~\eqref{listtheo1.1.3} on the decay of the linear semigroup may seem strong. They are satisfied in the cases where the set of trapped geodesics, the ones which do not meet the support of the damping, is small and hyperbolic in some sense. Several geometries satisfying~\eqref{listintro1.5} and~\eqref{listtheo1.1.3} have been studied in the literature, see the concrete examples of Figure~\ref{fig-1} and the references therein. Notice in particular that the example of domain with holes is particularly relevant for applications where we want to stabilize a nonlinear material with holes by adding a damping or a control in the external part. There is a huge literature about the damped wave equation and the purpose of the examples presented here is mainly to illustrate our theorem to non specialists. Moreover, the subject is growing quickly, giving more and more examples of geometries where we understand the effect of the damping and where we may be able to apply our results. We do not pretend to exhaustivity and refer to the bibliography of the more recent~\cite{Lin} for instance.

In some cases, the unique continuation property required in Proposition~\ref{prop-UCP-implique-th} can be obtained without considering analytic nonlinearities or conditions on the growth of $f$ as Hypothesis c) of Theorem~\ref{th1}. Instead, we require a particular geometry, which will be introduced more precisely in Section~\ref{section-UCP}.
\goodbreak

\begin{theo}\label{th2}
Consider the damped wave equation~\eqref{eq} in the framework of Assumptions~\eqref{listintro1.1}-\eqref{listintro1.5}. Assume in addition that:
\begin{enumerate}\alphenumi
\item \label{listtheo1.2.1}the function $(x,u)\mapsto f(x,u)$ is of class $\Cc^1\left(\overline{\Omega}\times \RR,\RR\right)$.
\item \label{listtheo1.2.2}there exists a pseudo-convex foliation of $\Omega$ in the sense of Definitions~\ref{defi-foliation1} or~\ref{defi-foliation2}.
\end{enumerate}
Then, the conclusions of Theorem~\ref{th1} hold.
\end{theo}

This result can be applied in several situations of Figure~\ref{fig-1}: the ``disk with two holes'', the ``peanut of rotation'' and the ``open book''. In these cases, the stabilization holds for any natural nonlinearity.

We expect that the decay rate $h_{R,\sigma}(t)$ is related to the linear decay rate $h(t)$ of Assumption~\eqref{listintro1.5}. We are able to obtain this link for the typical decays of the examples of Figure~\ref{fig-1}.

\begin{prop}\label{prop1}
Consider a situation where the stabilization stated in Theorems~\ref{th1} or~\ref{th2} holds. Then,
\begin{itemize}
\item if the decay rate of Assumption~\eqref{listintro1.5} satisfies $h(t)=\Oc({t^{-\alpha}})$ with $\alpha>1$, then the nonlinear equation admits a decay of the type $h_{R,\sigma}(t)= \Oc (t^{-\sigma\alpha})$.
\item if the decay rate of Assumption~\eqref{listintro1.5} satisfies $h(t)=\Oc \left(e^{-a {t^{1/\beta}}}\right)$ with $a>0$ and $\beta>0$, then the nonlinear equation admits a decay of the type $h_{R,\sigma}(t)\linebreak= \Oc \left(e^{-b\sigma t^{1/(\beta+1)}}\right)$ for some $b>0$.
\end{itemize}
\end{prop}

Notice that this result is purely local in the sense that the decay rate is obtained when the solution is close enough to $0$. Our proofs do not provide an explicit estimate of the time needed to enter this small neighborhood of $0$. Also notice that the loss in the power of the second case of Proposition~\ref{prop1} is due to an abstract setting: in the concrete examples, we may avoid this loss, see the remark below Lemma~\ref{lemma-appli2} and the concrete applications to the examples of Figure~\ref{fig-1}.

